%======================================================================
\section{Discussion}
\label{sec:discussion}
%======================================================================

\paragraph{Mechanism, not formatting alone.}
AEIC is useful here because it creates contractible observables, not because four labeled fields guarantee good reasoning. Static AEIC changes behavior when agents see a rule tied to those observables; Outcome Only shows that generic reward language is insufficient. At the same time, the signal ablations do not justify claiming that every field is independently indispensable.

\paragraph{Why the negative transfer matters.}
The GLM result indicates that prompt-mediated costs and credits are interpreted through a backend-specific policy. The contract's nominal numbers are therefore not yet a common currency. A credible next step is a calibration stage that estimates each backend's response surface on a development set before testing a fixed contract on held-out tasks. Mixed-model panels should be evaluated only after this calibration problem is addressed.

\paragraph{Relation to topology and deliberation.}
Production incentives and communication topology solve different problems. The former governs whether informative work is produced; the latter governs who can observe it and at what cost. The controlled information experiments begin to connect the two without turning topology into the paper's main causal variable. Future work can optimize information flow conditional on a calibrated reporting mechanism.

\paragraph{Limitations.}
The primary effect is established for one backend, three binary-answer QA domains, one prompt-mediated action menu, and majority aggregation. Declared actions are behavior generated by an LLM under instructions, not evidence of human-like preferences. Credits are bookkeeping quantities rather than monetary transfers, and process completeness is not factual verification. The confirmatory experiment tests a fixed static rule, not adaptive Stackelberg optimality or a general equilibrium. Finally, additional API calls and longer reports incur real compute cost even when token count is used only as a proxy.

\paragraph{Behavior before accuracy.}
The primary mechanism question is whether the contract changes follower action; collective correctness is a downstream outcome of the joint contract and execution path. The defensible claim is not that four AEIC labels inherently improve reasoning, but that an observable, enforceable, ex-ante process rule induces different resource behavior on the tested backend and is accompanied by a costly quality improvement.

\paragraph{Cost effectiveness.}
Static AEIC consumes 145.5 additional tokens per agent-round and yields 64 net question corrections relative to no mechanism. This accounting enables risk-sensitive deployment: low-risk items can retain cheap paths, while disputed or high-risk items receive verification budget. We have not solved the optimal routing threshold, so these figures expose the trade-off rather than establish optimal efficiency.

\paragraph{Broader implications.}
The framework is most appropriate when intermediate work can be made locally inspectable and when the executor can enforce meaningful differences between actions. It should not be used to manufacture confidence from structurally complete but unverifiable reports. Logging the offer, action, execution path, and settlement is therefore part of the contribution, not merely an engineering detail.
